\section{Removing the soft penalties up to physical time one}
\label{spb-section}

Work in dimension four, in current counting coordinates, on the chosen
infinite structural system. Flat spatial bounds also pass to explicit periodizations by summing the exponentially decaying lifted kernels. This is not a linear-weight similarity on a finite torus. The affine coefficient itself is taken on the infinite structural system, not as a literal periodic magnetic flux; arbitrary actual finite winding kernels are not identified with those periodizations.
Set $\mu=K^{-1}$. The affine background below means only the first
two formal polynomial jets of
\[
 U_\alpha(t;x)=\exp(i t x_\nu H/K^2),\qquad \alpha\ne\nu.
\]
No nonzero unbounded field is inserted in a common analytic ball.

\subsection{Actual families and the required inputs}
\label{spb-inputs}

Use the actual effective classical Hessians and moving-reference maps:
\begin{align}
 B_s(A)&=D_A^TD_A,&
 B_v(A)&=S_M''(e^{i\mu A})+D_AD_A^T,\notag\\
 T_s(A)&=B_s(A)+\mu^2J_AJ_A^T,&
 \widetilde V(A)&=B_v(A)+\bar C_A^TF_A\bar C_A,\notag\\
 F_A&=aI-S_{MK}''(M_K(e^{i\mu A})).
 \label{spb-families}
\end{align}
All Hessians are left-relative, including their coordinate-current
terms. The auxiliary vector $\widetilde V$ has the already established
first two affine jets of the actual soft vector; equality off shell
is not asserted.

For either pair $(W,B)=(T_s,B_s)$ or
$(W,B)=(\widetilde V,B_v)$, put $X_E=K^2E$ and write
\begin{equation}
 X_W=X_B+C_A^TF_AC_A,
 \qquad
 (C_A,F_A)=
 \begin{cases}
 (J_A^T,I),&\text{scalar},\\
 (K\bar C_A,\ aI-S_{MK}''(M_K(e^{i\mu A}))),&\text{vector}.
 \end{cases}
 \label{spb-coarse-factor}
\end{equation}
The output fiber of $C_A$ has a fixed color/site or color/bond
dimension per final block. Choose the fixed $a\ge R_q+1$, so that
$F_0\ge I$ in the vector case.

The proof uses the following established inputs, with constants
uniform in $M,K$:

\begin{enumerate}
\item The actual flat $q_M$ has an exponential current-coordinate
kernel, analytic double zero, Maxwell lower bound and uniform second
moments. Hence, with the componentwise physical gradient
$\nabla=K d_0$ and $H_0=I+\nabla^T\nabla$,
\begin{equation}
 cH_0\le I+X_{B,0}\le CH_0.
 \label{spb-base-energy}
\end{equation}
\item The repaired covariant-Hessian theorem gives first-two physical
$C^1$ vertices, with complete smooth source collars and full rooted
operator bounds. The scaled gauge-generator products have the same
property.
\item The physical maps $K\bar C_A$, $J_A$ and their first-two
background derivatives have uniform whole-block operator bounds and
fixed physical range. The actual coarse Hessian has uniform rooted
source kernels. Consequently $C_A^TF_AC_A$ has uniformly bounded
zero-order physical vertices and full rooted source bounds.
\end{enumerate}
These inputs concern actual effective actions and corrected tangents.
No bare Wilson replacement, full uniform nonlinear physical domain,
or arbitrary high-frequency bounded-source estimate is used.

\begin{theorem}
\label{spb-main}
Under these inputs, if $P_0$ is the appropriate fine scalar or bond
projector onto final block zero, then
\begin{equation}
 \left|[t^2]\operatorname{Tr}
 P_0\bigl(e^{-uX_W(t)}-e^{-uX_B(t)}\bigr)\right|
       \le Cu,\qquad 0<u\le1,
 \label{spb-physical-bound}
\end{equation}
uniformly in old depth $M$ and endpoint $K$. Equivalently,
\begin{equation}
 \left|[t^2]\operatorname{Tr}
 P_0\bigl(e^{-sW(t)}-e^{-sB(t)}\bigr)\right|
       \le C\frac{s}{K^2},\qquad 0<s\le K^2.
 \label{spb-fine-bound}
\end{equation}
In particular the absolute $ds/s$ integral of the difference on this
entire band is bounded uniformly.
\end{theorem}

\subsection{Shifted flat resolvents with physical spatial weights}
\label{spb-resolvents}

At zero the penalty in \eqref{spb-coarse-factor} is nonnegative and
bounded uniformly. Thus \eqref{spb-base-energy} holds for $W$ too.
Fix $E=B$ or $W$ and set
\[
 \mathcal B_0=I+X_{E,0},\qquad R_z=(\mathcal B_0-z)^{-1}.
\]
Use a sectorial contour with rays $\arg z=\pm\pi/4$, $|z|\ge1/4$,
joined by the circular arc of radius $1/4$ to the left of the positive
spectrum. The spectral theorem and energy comparison give
\begin{align}
 \|R_z\|&\le C(1+|z|)^{-1},\notag\\
 \|\nabla R_z\|+\|R_z\nabla^T\|
      &\le C(1+|z|)^{-1/2},\notag\\
 \|\nabla R_z\nabla^T\|&\le C.
 \label{spb-energy-resolvent}
\end{align}
For example,
$\|(\mathcal B_0+|z|)^{1/2}R_z
(\mathcal B_0+|z|)^{1/2}\|\le C$ on this contour.

Only linear physical weights are needed for the nonlocal flat kernel.
Conjugate by $e^{\theta\cdot x/K}$, with sufficiently small real
$\theta$. The analytic double zero gives the fine-symbol estimate
\[
 \|q_M(p+i\mu\theta)-q_M(p)\|
 \le C\{\mu|\theta|\sqrt{\delta(p)}+\mu^2|\theta|^2\}.
\]
After multiplication by $K^2$, this is a small form perturbation
relative to $H_0$. The local gradient term has the same estimate.
The coarse correction also does: its physical legs have bounded norm
and fixed range, and its middle flat kernel has a uniform coarse
analytic strip. The estimate remains true relative to $H_0+|z|$.
A weighted resolvent Neumann series therefore proves
\eqref{spb-energy-resolvent} uniformly after conjugation.
Commuting a weight through $\nabla$ adds only $O(|\theta|)$ and
is covered by the scalar resolvent estimate.

For each pair of blocks choose the signs of $\theta$ along their
separation. The restricted versions of all three bounds then acquire
$e^{-m d(i,j)}$, for a common $m>0$. Internal block diameters change
only the constant. These are whole-block Hilbert operator estimates,
with no fine-entry trace or growing-fiber count. At fixed $K$ the
operators are bounded lattice operators; finite-support approximation
justifies the weighted kernels. All constants are uniform in $K$.

\subsection{At most one unbounded derivative per vertex}
\label{spb-vertices}

Let $X_j=[t^j]X_E(tA)$, $j=1,2$, for a complete physical $C^1$
direction of norm at most $a_*$. The repaired Hessian decomposition,
multiplied by $K^2$, gives
\begin{equation}
 X_j=\nabla^T\overline Z_j+\overline Z_j^T\nabla
                  +\overline E_j,
 \qquad
 \|\overline Z_j\|+\|\overline E_j\|\le C_j a_*^j.
 \label{spb-first-order}
\end{equation}
Indeed $\overline Z_j=KZ_j$ and $\overline E_j=K^2E_j$;
the improved $\overline Z_2=O(K^{-1})$ is not needed.
The scaled gauge-generator terms have this form, and the physical
coarse penalty contributes bounded zero-order terms. Mixed second
coefficients follow by polarization.

For directions in smooth physical collars $x,y$, each factor retains
an exponential connecting-tree bound including its operator endpoints
and every source. The significant fact is that every summand in
\eqref{spb-first-order} has at most one unbounded physical derivative.
An arbitrary second-order form bound would not give the following
integrable resolvent powers.

Only two jets are needed. If a full uniform nonlinear domain has not
been provided, use the polynomial family
$X(t)=X_0+tX_1+t^2X_2$, or its mixed quadratic version.
Equation \eqref{spb-first-order} gives a small common parameter disk
on which $I+X(t)$ is sectorial with the chosen contour. At each fixed
$K$ this is ordinary holomorphic functional calculus for bounded
operators. Its first-two semigroup jets equal the actual ones because
only $X_0,X_1,X_2$ enter their resolvent derivatives. Thus this device
adds no nonlinear-domain hypothesis.

\subsection{Rooted semigroup derivatives}
\label{spb-semigroup}

Differentiated resolvents satisfy
\begin{align}
 \delta_xR_z&=-R_zX_xR_z,\notag\\
 \delta_x\delta_yR_z
  &=R_zX_xR_zX_yR_z+R_zX_yR_zX_xR_z-R_zX_{xy}R_z.
 \label{spb-resolvent-derivatives}
\end{align}
One first-order insertion in \eqref{spb-first-order} costs
$(1+|z|)^{-3/2}$ by \eqref{spb-energy-resolvent}; its zero-order
part costs $(1+|z|)^{-2}$. Two first-order insertions cost at worst
$(1+|z|)^{-2}$. If both derivatives face the middle resolvent, that
piece costs $O(1)$ and both exterior resolvents cost
$(1+|z|)^{-1}$. In the other arrangements there are two half-power
energy legs and one full resolvent. The single $X_{xy}$ insertion
again has exponent $-3/2$.

Insert block projections between the already defined factors
$R_z\nabla^T$, $\nabla R_z$ and the bounded vertex factors.
No sharp block cutoff is commuted through a derivative. The weighted
resolvent and vertex bounds then give, after an exponent reserve,
\begin{align}
 \|\delta_x(R_z)_{ij}\|
  &\le C(1+|z|)^{-3/2}e^{-m\tau(i,j,x)},\notag\\
 \|\delta_x\delta_y(R_z)_{ij}\|
  &\le C\{(1+|z|)^{-3/2}+(1+|z|)^{-2}\}
                         e^{-m\tau(i,j,x,y)}.
 \label{spb-rooted-resolvent}
\end{align}
Every intermediate sum is over block operator norms; a smaller
exponent pays the finite convolution multiplicities.

Integrate these formulas on the Dunford contour for
$e^{-u\mathcal B_0}$. The rays have positive real part, the joining
arc has fixed length, and both powers in
\eqref{spb-rooted-resolvent} are integrable at infinity, uniformly
for $0<u\le1$. Since $e^{-uX}=e^u e^{-u(I+X)}$, one obtains
\begin{align}
 \|\delta_x(e^{-uX})_{ij}\|
     &\le Ce^{-m'\tau(i,j,x)},\notag\\
 \|\delta_x\delta_y(e^{-uX})_{ij}\|
     &\le Ce^{-m'\tau(i,j,x,y)}.
 \label{spb-rooted-semigroup}
\end{align}
For instance the ray integrals of $(1+\lambda)^{-3/2}$ and
$(1+\lambda)^{-2}$ are bounded by fixed constants even when the
factor $e^{-u\lambda/\sqrt2}$ is discarded.

The undifferentiated semigroup has a uniform exponential restricted
block bound too. The small linear-weight form perturbation retains
the positive gradient energy, leaving numerical range bounded below
by $-C$. Its real-time norm is at most $e^{Cu}$, and the weight gives
the block decay. This uses only the flat generator.

This resolvent argument avoids the nonintegrable middle-time factor
that an overly crude split of a second Duhamel insertion would produce.
It also retains the source roots which a global Cauchy bound alone
would lose.

\subsection{Affine reconstruction and the coarse trace}
\label{spb-trace}

Use a fixed smooth physical partition for the affine potential.
Its direction in collar $x$ has physical $C^1$ norm at most
$C(1+|x|)$. Summing its first and mixed second directions against
\eqref{spb-rooted-semigroup} gives, for $n=0,1,2$,
\begin{equation}
 \|([t^n]e^{-uX(t)})_{ij}\|
 \le C(1+|i|+|j|)^p e^{-m''d(i,j)},\qquad 0<u\le1.
 \label{spb-affine-semigroup}
\end{equation}
At $u=0$ use the identity and zero differentiated coefficients.
The sums are absolutely convergent polynomial-jet sums, not evaluations
at a nonzero unbounded field.

The physical $C_A,F_A$ in \eqref{spb-coarse-factor} have the same
polynomial block bounds; $C_A$ has fixed physical range and order-one
norm. In the vector case ordinary averaging sends the affine connection
to the coarse affine connection, so the actual coarse Hessian supplies
the required polynomial/exponential $F$ kernels.

Retain $P_0$ in the exact physical-time Duhamel trace identity:
\begin{equation}
 \operatorname{Tr}_fP_0(e^{-uX_W}-e^{-uX_B})
  =-\int_0^u\operatorname{Tr}_c
       [FCe^{-vX_B}P_0e^{-(u-v)X_W}C^T]\,dv.
 \label{spb-coarse-duhamel}
\end{equation}
The trace on the right sums all coarse diagonal blocks.
Its degree-two coefficient has finitely many ordered terms.
Equation \eqref{spb-affine-semigroup} and the $C,F$ estimates connect
both coarse endpoints to the retained root block. Their absolute
sum is bounded by a fixed multiple of
\[
 \sum_{Y,Z}(1+|Y|+|Z|)^p
            e^{-m'''(|Y|+|Z|+d(Y,Z))}<\infty.
\]
Each diagonal trace is in a fixed-dimensional coarse output fiber.
A product through the fine block-zero Hilbert space costs that coarse
rank, not its growing fine dimension. Thus the integrand coefficient
is bounded uniformly in $u,v,M,K$, including the endpoints $v=0,u$.
The integral has length $u$, proving \eqref{spb-physical-bound}.

For compact source truncations the cyclic trace step is ordinary
finite-rank cyclicity with $P_0$ retained. The displayed absolute
rooted sum gives the infinite limit and permits coefficient extraction.
This completes the proof of Theorem~\ref{spb-main}.

\subsection{Consequence and remaining comparison}
\label{spb-scope}

Equation \eqref{spb-fine-bound} implies
\begin{equation}
 \int_0^{K^2}\left|[t^2]\operatorname{Tr}
 P_0(e^{-sW(t)}-e^{-sB(t)})\right|\frac{ds}{s}\le C.
 \label{spb-integrated}
\end{equation}
Together with the coarse determinant reductions, this removes the
soft penalties from the fine-to-physical-unit cutoff heat coefficient
up to a uniform finite part. The vector conclusion is actual through
the established affine matching; the scalar family is actual directly.

The remaining middle-band comparison concerns the translation-invariant
actual base $S_M''+DD^T$ and the ordinary scalar Laplacian.
This section does not replace that effective Hessian by a fixed bare
reference, compute its logarithmic coefficient, or prove the separate
physical-time tail beyond one. A theorem for a fixed bare integrand
is not imported as a uniform statement in the present family.

